\documentclass[11pt]{article}

\usepackage[margin=1in]{geometry}
\usepackage{amsmath, amssymb, amsthm}
\usepackage{mathtools}
\usepackage{enumitem}
\usepackage{hyperref}
\usepackage{graphicx}
\usepackage{tikz}
\usetikzlibrary{arrows.meta, positioning, calc}

\newcommand{\HS}{H^\star}

\title{Critical Summary: Decentralized Detection with Many Sensors}
\author{Oren Ram \\ Ziv Abraham}
\date{\today}

\begin{document}
\maketitle
\newpage

\section{Introduction}
\subsection{Background - Remote Hypothesis Testing}
- Add From Lecture -
\subsection{2}
In our work, we are going to discuss a problem of Decentrlized Hypothesis Testing, where many sensors observe an environment, compress the observation into an action, and send the actions to a fusion center (FC). The FC decides between two hypothesis - $H_1$ or $H_2$.\\
The paper we are going to discuss is Decentrlized Detection with Many Sensors: Optimality of Exchangeable and Identical Encoding Policies [1], and specifically, we are attending the finite number of sensors problem discussed in the paper.



\section{Problem Setup}

We ocnsider a problem of Binary Detection, as described in [1]. Assuming there are two hypothesis, $H_1$ and $H_2$, 
the true one can be regarded as $\HS$, a random varible with $P(\HS = H_1) = p$, $P(\HS = H_2) = 1-p$. The setup being tested:
\begin{itemize}
    \item $N < \infty$ sensors, each makes an observation of the environment deonted by $y^i \in \mathcal{Y}^i$ as the observation of the $i$-th sensor for $i \in \{1,\dots, N\}$.
    \item Each sensor quantizes its observation using a mapping $\gamma^i: \mathcal{Y}^i \to \mathcal{U}^i$, $\mathcal{U}^i$ being the set of actions the sensor can send.
    $\gamma^i$ is the desicion rule / policy of the $i$-th sensor. 
    \item The FC's policy is noted as $\gamma^0: \mathcal{U}^0 \times \mathcal{U}^1 \times \dots \mathcal{U}^N \to \{1,2\}$.\\
    Essentialy, the fusion center decides $u_0 = \gamma^0(u^{1:N}), u_0 \in \{1,2\}$.
    \item The cost function we will focus on is:
    \begin{align}
        c(H_j,u^0)= 
                        \begin{cases}
                            1 & \text{if } u^0 \neq j \\
                            0 & \text{if } u^0 = j
                        \end{cases}
    \end{align}
    \item Our goal is to minimize the expected cost:
    \begin{align}
        J^N(\gamma^{0:N}) \triangleq \mathbb{E}[c(H,u^0)]
    \end{align}
\end{itemize}
As the number of sensors $N$ increases, according to [1] and [2], $J^N(\gamma^{0:N})$ is bound by a decaying exponent,
regardless of the policies being used. This decay, and percisely how strong this decay is, is what carries the relevant information about
about the performence of a vector of policies, so the problem of minimizing the expected cost becomes maximizing the normlized error exponent - 
\begin{align}
    J^N_{EE}\bigg(\frac{\log(J^N(\gamma^{0:N}))}{N}\bigg)
\end{align}
\newpage
\subsection{Assumptions}
In order to make the problem tractable, the authors of the paper consider a few assumptions:
\begin{enumerate}[label=(\roman*)]
    \item $y^1, \dots, y^N$ are i.i.d conditioned on $H$
    \item Under either hypothesis, sensor $i$'s observation has a probability density, measured against one common reference measure shared across both hypothesis.
    \item $\mathcal{U}^i = \mathcal{U}$ for $i = 1,\dots ,N$, and $\mathcal{U}$ is a finite set.
    \item The likelihood ratio:
    \begin{align}
        l^i = L(y^i) := \frac{f_{y^i|H}(y^i|H_2)}{f_{y^i|H}(y^i|H_1)} \text{a.s}
    \end{align}
    Is well behaving and invertible.
\end{enumerate}
While (i) is a common assumption in order to simplify the math and making it tractable, assumptions (ii) and (iv) are essential in order for $l^i$ to be a sufficient statistic of $y^i$, as will be discussed later.
Assumption (iii) dictates that all sensors have the same action set, underlying assumption about the compression and rate being used. This will be discussed furthermore in section *, where we will offer some modifications for the solution setup.

\section{Sufficient Statistics, Existence of an Optimal Encoding Policy}

\subsection{Exponential Decay of the Error Probability}

\subsection{$l^i$ is a Sufficient Statistic}

$l^i$ is a sufficient statistic of $H$: once its value is known, nothing else
about the raw observation $y^i$ can further move one's belief about which
hypothesis is true. With only two hypotheses on the table, the entire way in which
the observation's density can vary with $H$ is exhausted by a single comparison -
how large the density is under $H_2$ relative to how large it is under $H_1$, at
that particular observation. A single number is therefore inherently enough to
hold everything that changes with $H$, and that number is, by construction,
exactly $l^i$. (How this shows up concretely inside the paper's own machinery -
via Lemma~1 -- is worth treating separately once we get there.)

\subsubsection*{No information loss: a Markov-chain and DPI reading}
Sufficiency means that any two raw observations producing the same $l^i$ are, for
the purpose of distinguishing $H_1$ from $H_2$, interchangeable: whatever detail of
$y^i$ is discarded in forming $l^i$ carries no evidential weight about $H$. This
has a clean information-theoretic statement. Since $l^i$ is a deterministic
function of $y^i$, the chain
\[
    H \;\to\; y^i \;\to\; l^i
\]
is trivially Markov, and the Data Processing Inequality (DPI) gives
\begin{equation}
    I(H; l^i) \le I(H; y^i).
\end{equation}
Equality holds \emph{iff} $l^i$ is sufficient for $H$ -- equivalently, iff
$H \to l^i \to y^i$ is \emph{also} a Markov chain, i.e.\ $y^i$ carries no further
information about $H$ once $l^i$ is fixed, gives equality
\begin{equation}
    I(H; l^i) = I(H; y^i).
\end{equation}

\subsection{Optimality of Threshold Policies and the MAP Decoder}
\subsubsection*{Lemma 1}
Let Assumption 1 hold. Then, for any fusion rule $\gamma^0 \in \Gamma^0$,
\begin{equation}
    \inf_{\gamma^{1:N} \in \Gamma^N} J^N(\gamma^0, \gamma^{1:N})
    = \inf_{\gamma^{1:N} \in \Gamma^N_{TS}} J^N(\gamma^0, \gamma^{1:N}).
\end{equation}
Where $\Gamma^N_{TS}$ is the class of threshold policies.
Fixing every sensor's policy except sensor $i$'s, sensor $i$'s optimization is:

\begin{align}
    \inf_{u^i \in \mathcal U} \bigg\{\sum_{j=1,2} P(H_j \mid y^i)\, E_{\gamma^{-i\star}}\big[c(H_j, \gamma^0(u^{1:N})) \mid H_j\big]\bigg\}
\end{align}

By Bayes' rule,
$P(H_j \mid y^i)$ depends on $y^i$ only through the likelihood ratio $l^i$ --
this is exactly the sufficiency of $l^i$ from Section [X] showing up again,
now inside the optimization itself. Substituting reduces the problem to
\begin{equation}
    \inf_{u^i \in \mathcal U} \left[ g^i(H_1, u^i) + g^i(H_2, u^i)\, l^i\,
    \frac{P(H_2)}{P(H_1)} \right],
\end{equation}
where $g^i(H_j, u^i)$ is the expected cost of action $u^i$ under $H_j$ given the
other sensors' fixed policies. This is linear in $l^i$ for each action, so the
best action is whichever line is lowest at that value of $l^i$, switching only
at the points where lines cross -- exactly a threshold rule.
\subsubsection*{Lemma 2}

\subsection{Existence of an Optimal Solution}
Lemmas 1 and 2 tell us where to look; Theorem 1, via the compactness result of Lemma 3, is what guarantees something is actually there to find

\section{Considering Rate}



\end{document}
